\section{Introduction}\label{sec:introduction}

Particle-in-cell (PIC) simulations are a useful numerical method for the study of plasma dynamics since they can record particle positions, momenta, and electromagnetic fields, which are sampled along particle trajectories, thus generating large sets of multivariate time series. As the number of tracked particles, recorded time steps, and available variables increases, it becomes insufficient to examine individual trajectories or to compare a small number of two-dimensional projections for systematic analysis. For this reason, automated feature extraction and the grouping of particles with similar trajectory characteristics constitute an important aspect of PIC data analysis.

Unsupervised learning allows this task to be carried out without the need for previously defined particle labels. By constructing suitable data representations and grouping samples by similarity, large trajectory datasets can be organized into more manageable subsets. Markidis et al. introduced an automated workflow for classifying particle trajectories from plasma simulations by combining the fast Fourier transform, principal component analysis, K-means clustering, and silhouette analysis. This method was applied to electron trajectories in a magnetic reconnection simulation, which laid the foundation for the investigation of particle trajectory datasets using dimensionality reduction and clustering~\cite{markidis2020}.

Deep representation learning offers another approach to processing complex trajectory data. Autoencoders learn low-dimensional features through an input reconstruction task, so that high-dimensional samples can be represented in a compact space. Deep Embedded Clustering (DEC) combines representation learning with cluster assignment, iteratively modifying both the feature representation and the positions of the cluster centers in order to achieve the clustering aim. This method provides a framework in which the representation used to group samples is learned from the data~\cite{xie2016}.

The present study is based on the internship report \textit{Exploring Particle In Cell Simulation Data Using Artificial Intelligence} by Medina Moreira. That work looked into unsupervised particle trajectory analysis using position and kinetic energy time series, spectrogram representations, convolutional autoencoders, and deep clustering. It also pointed out that there was a lack of information regarding momentum and the electromagnetic fields in the dataset and proposed that the workflow be extended by including further variables as input channels. This suggestion from the report serves as the foundation for the present study, which applies the existing framework to more comprehensive descriptions of particle trajectories and investigates the related data processing and clustering performance~\cite{moreiraReport}.

Extending the input to additional variables leads to further data-processing challenges. Since position, momentum, and electromagnetic fields have different numerical scales and temporal characteristics, their relative contributions are sensitive to preprocessing. The inclusion of more variables also raises the dimensionality of each sample and may result in correlated or redundant information. Therefore, channel normalization, feature extraction, and clustering evaluation must be considered together when constructing a multivariate trajectory representation. Following the previous workflow, this study uses spectrograms to represent the evolution of local frequency content within each time series. The spectrograms of different variables are combined into a multichannel input for feature learning with a convolutional autoencoder.

The input variables considered here comprise particle position \(x\), momentum components \(p_x\) and \(p_y\), electric field components \(E_x\) and \(E_y\), and the magnetic field component \(B_z\), with the fields recorded along particle trajectories. The workflow begins by normalizing the variables and then transforms the time-series data into spectrograms. A convolutional autoencoder is then used to learn low-dimensional representations, which are used to initialize cluster centers through K-means. Subsequent deep clustering refinement retains a reconstruction constraint while updating both the representation and cluster assignments. The spectrogram-based workflow determines the number of clusters using an elbow criterion. This selected value is then applied to the reference methods so that the number of clusters remains consistent across the comparisons.

\begin{figure}[htbp]
\centering
\figureasset[\linewidth]{figures/workflow-no-trajectories.png}{Overview of the multivariate PIC deep clustering workflow}
\caption[Overview of the multivariate PIC deep clustering workflow]{Overview of the spectrogram-based deep clustering workflow for multivariate PIC particle trajectories. Time series of particle position, momentum components, and local electromagnetic fields are normalized and transformed into multichannel log-spectrograms. A convolutional autoencoder learns latent representations, from which the number of clusters is selected using an elbow criterion. K-means initialization is followed by deep clustering refinement with a reconstruction constraint to obtain the final cluster assignments. Clustering quality is evaluated primarily using the silhouette score, with reconstruction errors, visualizations, and cluster statistics providing supplementary diagnostics. All plots are schematic.}\label{fig:workflow}
\end{figure}

Because predefined particle labels are unavailable, clustering quality is mainly assessed through internal validation. The silhouette score is used as the main evaluation metric. This score measures both cohesion within clusters and separation between clusters by comparing, for each sample, its distances to samples in its own cluster and to those in its nearest neighboring cluster~\cite{rousseeuw1987}. Its interpretation depends on the feature space and distance metric used for evaluation. Supplementary diagnostic information is obtained from reconstruction errors, cluster-size distributions, and dimensionality-reduction visualizations. As a reference approach to help assess the effect of the spectrogram representation, a one-dimensional convolutional autoencoder that operates directly on the original time series is included. A further linear baseline is provided by PCA followed by K-means using the same spectrogram inputs as the main method.

This study aims to apply the earlier spectrogram-based deep clustering procedure to multivariate particle data and to evaluate its capabilities in terms of feature extraction and clustering. The rest of the report is structured as follows. Section~\ref{sec:data} deals with the dataset, input variables, and preprocessing methods. Section~\ref{sec:methodology} presents the construction of multichannel spectrograms, the convolutional autoencoder, the deep clustering procedure, and the strategy for selecting the number of clusters. Section~\ref{sec:experiments} sets out the experimental design and evaluation method. Section~\ref{sec:results} presents the clustering results together with additional diagnostic information. Section~\ref{sec:conclusions} contains the conclusions, limitations, and implications. Section~\ref{sec:ai-assistance} provides the declaration of AI assistance. Further details regarding the architecture and training settings are given in Appendix~\ref{app:implementation}, and the final cluster statistics are listed in Appendix~\ref{app:clusters}.
